\documentclass[draftproposal,form,masculine]{pncdi-pd}

% \pncdiidentifier{PN-IV-RU-SC-PD-2026-1}

\begin{document}

\textbf{Anexa 5.1 -- Declarație privind încadrarea în definiția organizației de cercetare}

\bigskip

\begin{center}
\textbf{Declarație privind încadrarea în definiția organizației de cercetare}
\end{center}

\gendered{Subsemnatul}{Subsemnata} \placeholder[15]{} \textit{(numele şi prenumele reprezentantului
legal al organizației de cercetare)}, în calitate de \placeholder[15]{} \textit{(funcția
reprezentantului legal)} al \placeholder[12]{} \textit{(denumirea completă a organizației
de cercetare)}, declar pe proprie răspundere că sunt îndeplinite următoarele condiții
privind încadrarea în definiția organizației de cercetare:
\begin{itemize}[leftmargin=*]
\item
Este o entitate\footnotemark[1], indiferent de statutul său juridic (organizație de drept public sau
privat sau de modalitatea de finanțare, al cărei obiectiv principal este de a efectua în
mod independent cercetare fundamentală, cercetare industrială sau dezvoltare experimentală,
sau de a disemina la scară largă rezultatele unor astfel de activități, prin predare,
publicare, sau transfer de cunoștințe.

\footnotetext[1]{Universități sau institute de cercetare, agenții de transfer de tehnologie,
intermediari pentru inovare, entități de colaborare fizice sau virtuale orientate spre
cercetare, inclusiv nucleele de cercetare constituite în cadrul unor spitale sau muzee, care
îndeplinesc criteriile cuprinse în definiția organizației de cercetare
}

\item
În cazul în care entitatea desfășoară și activități economice, finanțarea, costurile și
veniturile activităților economice respective trebuie să fie contabilizate separat.

\item
Întreprinderile care pot exercita o influență decisivă asupra unei astfel de entități,
de exemplu, în calitate de acționari sau asociați, nu pot beneficia de acces preferențial
la rezultatele generate de aceasta.
\end{itemize}

\bigskip

\textbf{Declarație pe proprie răspundere, sub sancțiunile aplicate faptei de fals în acte publice.}

\bigskip

{
\noindent \begin{tabular}{@{}p{0.45\textwidth}p{0.55\textwidth}@{}}
Data: & \textbf{\rodate} \\[1em]
Reprezentant legal & Funcția: \textbf{Funcția} \\[0.3em]
& Numele și prenumele: \textbf{Nume, Prenume} \\[0.3em]
& Semnătura \\[2em]
Director de proiect & Numele și prenumele: \textbf{Nume, Prenume} \\[0.3em]
& Semnătura
\end{tabular}
}

\end{document}

% kate: default-dictionary ro_RO;
